\documentclass[../../数学分析培优讲义(答案版).tex]{subfiles}

\begin{document}

\setcounter{chapter}{8}
\pagestyle{fancy}

\chapter{函数项级数、幂级数、Fourier 级数}

\section{一致收敛性}

\begin{example}
试判断下列函数列在指定区间上的一致收敛性:
\begin{enumerate}[label=(\arabic*)]
    \item $f_n(x)=x^n-x^{n+1}$, $x\in\lbrack 0,1\rbrack$;
    \item $f_n(x)=\dfrac{x}{1+nx^2}$, $x\in(-\infty,+\infty)$;
    \item $f_n(x)=n^\alpha e^{-nx}$, $x\in\lbrack 0,1\rbrack$.
\end{enumerate}
\end{example}

\begin{proof}[解]
\begin{enumerate}[label=(\arabic*)]
    \item 极限函数为零，且
    \[
        \max_{0\leq x\leq1}x^n(1-x)
        =\dfrac1{n+1}\left(\dfrac n{n+1}\right)^n\longrightarrow0,
    \]
    故一致收敛。
    \item 极限函数为零，且
    \[
        \sup_{x\in\mathbb R}\dfrac{|x|}{1+nx^2}=\dfrac1{2\sqrt n}\longrightarrow0,
    \]
    故一致收敛。
    \item 当 $\alpha<0$ 时，$0\leq f_n(x)\leq n^\alpha\to0$，故一致收敛于零。当 $\alpha=0$ 时，逐点极限在 $x=0$ 处为 $1$、在 $x>0$ 处为零，不连续，故不一致收敛；当 $\alpha>0$ 时，$f_n(0)=n^\alpha$，连逐点有限收敛也不成立。
\end{enumerate}
\end{proof}

\begin{example}
试判断下列函数列在指定区间上的一致收敛性:
\begin{enumerate}[label=(\arabic*)]
    \item $f_n(x)=\sin\dfrac{1+nx}{2n}$, $x\in(-\infty,+\infty)$;
    \item $f_n(x)=n\sin\dfrac{1}{nx}$, $x\in\lbrack 1,+\infty)$;
    \item $f_n(x)=\dfrac{\ln(nx)}{nx^2}$, $x\in(1,+\infty)$;
    \item $f_n(x)=\dfrac{n}{x}\ln\left(1+\dfrac{x}{n}\right)$, $x\in(0,10)$.
\end{enumerate}
\end{example}

\begin{proof}[解]
\begin{enumerate}[label=(\arabic*)]
    \item 极限为 $\sin(x/2)$，且由正弦函数的 Lipschitz 性，
    \[
        \sup_{x\in\mathbb R}\left|f_n(x)-\sin\dfrac x2\right|\leq\dfrac1{2n},
    \]
    故一致收敛。
    \item 极限为 $1/x$。由 $|\sin t-t|\leq|t|^3/6$，
    \[
        \sup_{x\geq1}\left|n\sin\dfrac1{nx}-\dfrac1x\right|
        \leq\dfrac1{6n^2},
    \]
    故一致收敛。
    \item 极限为零。充分大时 $x\mapsto\ln(nx)/x^2$ 在 $x>1$ 上递减，故
    \[
        \sup_{x>1}|f_n(x)|=\dfrac{\ln n}{n}\longrightarrow0.
    \]
    \item 极限为 $1$，且 $0<\ln(1+t)-t+t^2/2<t^2/2$ 给出
    \[
        \sup_{0<x<10}|f_n(x)-1|\leq\dfrac5n\longrightarrow0.
    \]
\end{enumerate}
\end{proof}

\begin{example}
证明下列函数列在指定区间上的非一致收敛性:
\begin{enumerate}[label=(\arabic*)]
    \item $f_n(x)=x^n-x^{2n}$, $x\in\lbrack 0,1\rbrack$;
    \item $f_n(x)=x^n+x^{2n}-2x^{3n}$, $x\in\lbrack 0,1\rbrack$;
    \item $f_n(x)=nx(1-x)^n$, $x\in\lbrack 0,1\rbrack$;
    \item $f_n(x)=n\left(\sqrt{x+\dfrac{1}{n}}-\sqrt{x}\right)$, $x\in(0,+\infty)$.
\end{enumerate}
\end{example}

\begin{proof}[解]
四列均逐点趋于零，但可分别选取
\[
    x_n=2^{-1/n},\qquad x_n=2^{-1/n},\qquad
    x_n=\dfrac1{n+1},\qquad x_n=\dfrac1{n^2}.
\]
此时前两列的函数值分别恒为 $1/4$ 与 $1/2$；第三列趋于 $e^{-1}$；第四列为
\[
    \dfrac1{\sqrt{n^{-2}+n^{-1}}+n^{-1}},
\]
甚至趋于 $+\infty$。因此都不一致收敛。
\end{proof}

\begin{example}
试判断下列级数 $\displaystyle I=\sum\limits_{n=1}^{\infty}f_n(x)$ 在指定区间上的一致收敛性:
\begin{enumerate}[label=(\arabic*)]
    \item $\displaystyle\sum\limits_{n=1}^{\infty}x^\alpha e^{-nx^2}$, $x\in(0,+\infty)$;
    \item $\displaystyle\sum\limits_{n=1}^{\infty}e^{-(x-n)^2}$, $x\in\lbrack a,b\rbrack$;
    \item $\displaystyle\sum\limits_{n=1}^{\infty}\dfrac{x^n}{4^n}\ln\left(1+\dfrac{x^2}{n^2}\right)$, $x\in\lbrack -4,4\rbrack$;
    \item $\displaystyle\sum\limits_{n=1}^{\infty}\dfrac{x}{n^p+n^q x^2}$, $x\in\lbrack 0,+\infty)$.
\end{enumerate}
\end{example}

\begin{proof}[解]
\begin{enumerate}[label=(\arabic*)]
    \item 当 $\alpha>0$ 时，
    \[
        \sup_{x>0}x^\alpha e^{-nx^2}=C_\alpha n^{-\alpha/2}.
    \]
    因此 $\alpha>2$ 时由 Weierstrass 判别一致收敛。若 $\alpha\leq2$，在尾项中取 $x=N^{-1/2}$ 并考察 $N\leq n\leq2N$，尾和不趋于零，故不一致收敛。
    \item 设 $C=\max\{|a|,|b|\}$。充分大时
    $e^{-(x-n)^2}\leq e^{-(n-C)^2}$，故一致收敛。
    \item 对 $|x|\leq4$，
    \[
        \left|\dfrac{x^n}{4^n}\ln\left(1+\dfrac{x^2}{n^2}\right)\right|
        \leq\dfrac{16}{n^2},
    \]
    故一致绝对收敛。
    \item
    \[
        \sup_{x\geq0}\dfrac{x}{n^p+n^qx^2}=\dfrac1{2n^{(p+q)/2}}.
    \]
    因而 $p+q>2$ 时一致收敛；反之取 $x=N^{(p-q)/2}$ 并对
    $N\leq n\leq2N$ 求和，尾和不趋于零。故恰在 $p+q>2$ 时一致收敛。
\end{enumerate}
\end{proof}

\begin{example}
试判断下列级数在指定区间上的一致收敛性:
\begin{enumerate}[label=(\arabic*)]
    \item $\displaystyle\sum\limits_{n=1}^{\infty}\dfrac{2nx}{1+n^\alpha x}$, $-\infty<x<+\infty$, $\alpha>\varphi$;
    \item $\displaystyle\sum\limits_{n=1}^{\infty}\dfrac{x\sin nx}{\sqrt{1+n^2(1+nx^2)}}$, $x\in(-\infty,+\infty)$;
    \item $\displaystyle\sum\limits_{n=1}^{\infty}\dfrac{x^2}{1+n^2x}\sin\dfrac{n}{x}$, $x\in(0,1)$;
    \item $\displaystyle\sum\limits_{n=1}^{\infty}\dfrac{(-1)^{n+1}}{2n+\sin x}$, $x\in(-\infty,+\infty)$.
\end{enumerate}
\end{example}

\begin{proof}[解]
\begin{enumerate}[label=(\arabic*)]
    \item 当前题面在 $x=-n^{-\alpha}$ 处分母为零，故函数甚至未在整个实轴上定义，无法讨论该区间上的一致收敛性；区间或分母符号需要核对。
    \item
    \[
        \left|\dfrac{x\sin nx}{\sqrt{1+n^2(1+nx^2)}}\right|
        \leq\dfrac{|x|}{n\sqrt{1+nx^2}}\leq\dfrac1{n^{3/2}},
    \]
    故一致绝对收敛。
    \item 对 $0<x<1$，
    \[
        \left|\dfrac{x^2}{1+n^2x}\sin\dfrac nx\right|\leq\dfrac{x}{n^2}\leq\dfrac1{n^2},
    \]
    故一致绝对收敛。
    \item $b_n(x)=1/(2n+\sin x)$ 关于 $n$ 递减，且
    $\sup_xb_n(x)\leq1/(2n-1)\to0$，故由一致 Leibniz 判别法一致收敛。
\end{enumerate}
\end{proof}

\begin{example}
试判断下列级数在指定区间上的非一致收敛性:
\begin{enumerate}[label=(\arabic*)]
    \item $\displaystyle\sum\limits_{n=1}^{\infty}\dfrac{x^3}{(1+x^3)^n}$, $x\in(0,+\infty)$;
    \item $\displaystyle\sum\limits_{n=1}^{\infty}\dfrac{x}{1+n^3x^3}$, $x\in(0,+\infty)$;
    \item $\displaystyle\sum\limits_{n=1}^{\infty}nx^2e^{-nx}$, $x\in(0,+\infty)$;
    \item $\displaystyle\sum\limits_{n=1}^{\infty}\dfrac{x}{x^2-nx+n^2}$, $x\in(1,+\infty)$.
\end{enumerate}
\end{example}

\begin{proof}[解]
\begin{enumerate}[label=(\arabic*)]
    \item 几何级数的余项为 $(1+x^3)^{-N}$，其在 $x>0$ 上的上确界为 $1$，故不一致收敛。
    \item 取 $x=1/N$，则从 $n=N$ 到 $2N$ 的尾和为
    \[
        \sum\limits_{n=N}^{2N}\dfrac{N^{-1}}{1+(n/N)^3}\geq c>0.
    \]
    \item 仍取 $x=1/N$，有
    \[
        \sum\limits_{n=N}^{2N}nx^2e^{-nx}geq c>0.
    \]
    \item 取 $x=N$，对 $N\leq n\leq2N$，每项均为 $N^{-1}$ 量级，故相应尾和有正下界。
\end{enumerate}
四者均违反一致 Cauchy 准则。
\end{proof}

\begin{example}
试用 Cauchy 收敛准则判别下列级数在指定区间上的非一致收敛性:
\begin{enumerate}[label=(\arabic*)]
    \item $\displaystyle\sum\limits_{n=1}^{\infty}\dfrac{n}{x}e^{-n^2/x}$, $x\in(0,+\infty)$;
    \item $\displaystyle\sum\limits_{n=1}^{\infty}\dfrac{\sin nx}{n}$, $x\in\lbrack 0,2\pi\rbrack$;
    \item $\displaystyle\sum\limits_{n=1}^{\infty}\dfrac{\sin nx}{e^{n^2x}}$, $x\in(0,+\infty)$;
    \item $\displaystyle\sum\limits_{n=1}^{\infty}\dfrac{\sqrt{\pi}}{n^2}\sin\dfrac{\pi}{n^2}$, $x\in(0,+\infty)$.
\end{enumerate}
\end{example}

\begin{proof}[解]
\begin{enumerate}[label=(\arabic*)]
    \item 取 $x=N^2$ 并令 $N\leq n\leq2N$，则尾和
    $\sum(n/N^2)e^{-n^2/N^2}$ 有固定正下界。
    \item 取 $x=1/N$，则 $N\leq n\leq2N$ 时 $nx\in[1,2]$，相应尾和有固定正下界。
    \item 取 $x=1/N^2$，利用 $\sin(nx)\sim nx$，从 $N$ 到 $2N$ 的尾和仍有固定正下界。
    \item 当前级数与 $x$ 无关，并且
    \[
        \left|\dfrac{\sqrt\pi}{n^2}\sin\dfrac\pi{n^2}\right|
        \leq\dfrac{\pi^{3/2}}{n^4}.
    \]
    因而它实际上一致绝对收敛，与题目要求“非一致收敛”矛盾，题面应有漏项。
\end{enumerate}
\end{proof}

\begin{example}
试用端点法判别下列级数在指定区间上的非一致收敛性:
\begin{enumerate}[label=(\arabic*)]
    \item $\displaystyle\sum\limits_{n=1}^{\infty}e^{-nx}\cos nx$, $x\in(0,1)$;
    \item $\displaystyle\sum\limits_{n=1}^{\infty}\dfrac{\ln(nx)}{1+n^3\ln^2x}$, $x\in\left(\dfrac{1}{2},1\right)$;
    \item $\displaystyle\sum\limits_{n=1}^{\infty}\dfrac{1}{(1+nx)^2}$, $x\in(0,1)$;
    \item $\displaystyle\sum\limits_{n=1}^{\infty}\left(1+n\dfrac{\sin x}{x}\right)^2$, $x\in(0,\pi)$.
\end{enumerate}
\end{example}

\begin{proof}[解]
\begin{enumerate}[label=(\arabic*)]
    \item 当 $x\to0^+$ 时，各项趋于 $1$，相应和函数无界，故不可能是一致极限。
    \item 对每个 $x<1$ 级数收敛，但 $x\to1^-$ 时第 $n$ 项趋于 $\ln n$，端点处连通项趋零条件都失败，故不一致收敛。
    \item 当 $x\to0^+$ 时，$\sum(1+nx)^{-2}$ 发散到无穷；而各部分和有界，故不可能一致收敛。
    \item 当前通项对每个 $x\in(0,\pi)$ 都不趋于零，原级数逐点发散，不能称为“非一致收敛”；指数的正负号需要核对。
\end{enumerate}
\end{proof}

\begin{example}
试用 Abel--Dirichlet 判别法判别下列级数在指定区间上的一致收敛性:
\begin{enumerate}[label=(\arabic*)]
    \item $\displaystyle\sum\limits_{n=1}^{\infty}\dfrac{(-1)^n(x^2+n)}{n^2}$, $x\in\lbrack a,b\rbrack$;
    \item $\displaystyle\sum\limits_{n=1}^{\infty}\dfrac{(-1)^{n+1}e^{-nx}}{\sqrt{n+x^2}}$, $x\in\lbrack 0,+\infty)$;
    \item $\displaystyle\sum\limits_{n=1}^{\infty}\dfrac{(-1)^n}{\sqrt n}\sin\left(1+\dfrac{x}{n}\right)$, $x\in\lbrack a,b\rbrack$;
    \item $\displaystyle\sum\limits_{n=1}^{\infty}\dfrac{\cos nx}{n^p}$, $x\in\lbrack \delta,2\pi-\delta\rbrack$, $p\in(0,1\rbrack$.
\end{enumerate}
\end{example}

\begin{proof}[解]
\begin{enumerate}[label=(\arabic*)]
    \item 分解为 $x^2\sum(-1)^n/n^2+\sum(-1)^n/n$；前者由 M 判别、后者与 $x$ 无关，故一致收敛。
    \item $e^{-nx}/\sqrt{n+x^2}$ 关于 $n$ 递减，且其关于 $x\geq0$ 的上确界为 $1/\sqrt n\to0$，故一致收敛。
    \item 写成
    \[
        \sin1\sum\limits_{n=1}^{\infty}\dfrac{(-1)^n}{\sqrt n}
        +\sum\limits_{n=1}^{\infty}\dfrac{(-1)^n}{\sqrt n}
        \left[\sin\left(1+\dfrac xn\right)-\sin1\right].
    \]
    第二级数在有限区间上由 $C/n^{3/2}$ 控制，故一致收敛。
    \item 三角和的部分和在 $[\delta,2\pi-\delta]$ 上一致有界，而 $n^{-p}\downarrow0$，由一致 Dirichlet 判别得结论。
\end{enumerate}
\end{proof}

\begin{example}
试由和函数的不连续性判别下列级数在指定区间上的非一致收敛性:
\begin{enumerate}[label=(\arabic*)]
    \item $\displaystyle S(x)=\sum\limits_{n=0}^{\infty}\dfrac{x^4}{(1+x^4)^n}$, $x\in\lbrack 0,1\rbrack$;
    \item $\displaystyle S(x)=\sum\limits_{n=0}^{\infty}x^n(1-x)$, $x\in\lbrack 0,1\rbrack$;
    \item $\displaystyle S(x)=\sum\limits_{n=1}^{\infty}x^n\ln x$, $x\in(0,1\rbrack$;
    \item $\displaystyle S(x)=\sum\limits_{n=2}^{\infty}\left(x^{\frac{1}{2n-1}}-x^{\frac{1}{2n-3}}\right)$, $x\in\lbrack -1,1\rbrack$.
\end{enumerate}
\end{example}

\begin{proof}[解]
\begin{enumerate}[label=(\arabic*)]
    \item 当 $x>0$ 时和为 $1+x^4$，而 $x=0$ 时和为 $0$，在零点不连续。
    \item 当 $0\leq x<1$ 时和为 $1$，而 $x=1$ 时和为 $0$，在右端点不连续。
    \item 当 $0<x<1$ 时和为 $x\ln x/(1-x)$，其在 $x\to1^-$ 时趋于 $-1$；但 $x=1$ 时每项为零，和为零。
    \item 前 $N$ 项裂项为 $x^{1/(2N-1)}-x$。极限在 $x>0$ 时为 $1-x$，在 $x=0$ 时为零，在 $x<0$ 时为 $-1-x$，故在零点不连续。
\end{enumerate}
各项均连续，而和函数不连续，所以均不一致收敛。
\end{proof}

\begin{example}
设 $f(x)$ 是 $(-\infty,+\infty)$ 上的连续函数,
\[
    f_n(x)=\sum\limits_{k=0}^{n-1}\dfrac{1}{n}f\left(x+\dfrac{k}{n}\right).
\]
证明:函数列 $\{f_n(x)\}$ 在任何有限区间上一致收敛.
\end{example}

\begin{proof}[解]
在任意有限区间 $[a,b]$ 上，$f$ 在 $[a,b+1]$ 上一致连续。$f_n(x)$ 是函数
$t\mapsto f(x+t)$ 在 $[0,1]$ 上等距分割的 Riemann 和，故
\[
    f_n(x)\rightrightarrows\int_0^1f(x+t)\,\mathrm{d}t
\]
且误差由 $f$ 在 $[a,b+1]$ 上的连续模一致控制。
\end{proof}

\begin{example}
设函数 $f(x)$ 在 $(-\infty,+\infty)$ 上有连续导函数 $f'(x)$,
\[
    f_n(x)=e^n\left[f(x+e^{-n})-f(x)\right].
\]
证明:$\{f_n(x)\}$ 在任一有限开区间 $(a,b)$ 内一致收敛于 $f'(x)$.
\end{example}

\begin{proof}[解]
写成
\[
    f_n(x)=e^n\int_0^{e^{-n}}f'(x+t)\,\mathrm{d}t.
\]
在包含 $[a,b]$ 的稍大紧区间上，$f'$ 一致连续，故
\[
    \sup_{a<x<b}|f_n(x)-f'(x)|
    \leq\sup_{\substack{a<x<b\\0\leq t\leq e^{-n}}}|f'(x+t)-f'(x)|\longrightarrow0.
\]
\end{proof}

\begin{example}[Heine 定理的推广]
试证:$x\to a$ 时 $f(x,y)\rightrightarrows\varphi(y)$ 的充要条件是对任意 $\{x_n\}\to a$,有 $f(x_n,y)\rightrightarrows\varphi(y)$ $(n\to\infty)$.
\end{example}

\begin{proof}[解]
必要性直接由定义得到。反之若不一致趋于 $\varphi$，则存在 $\varepsilon_0>0$，以及趋于 $a$ 的点列 $x_n$，使
\[
    \sup_y|f(x_n,y)-\varphi(y)|\geq\varepsilon_0.
\]
这与题设对任意点列均一致收敛矛盾，故条件充分。
\end{proof}

\begin{example}
若 $f_n(x)$ 在 $\lbrack a,b\rbrack$ 上可积,$f(x)$ 与 $g(x)$ 在 $\lbrack a,b\rbrack$ 上都可积,且
\[
    \lim\limits_{n\to\infty}\int_a^b|f_n(x)-f(x)|^2\,\mathrm{d}x=0.
\]
设
\[
    h(x)=\int_a^x f(t)g(t)\,\mathrm{d}t,
    \qquad
    h_n(x)=\int_a^x f_n(t)g(t)\,\mathrm{d}t.
\]
证明:$h_n(x)$ 在 $\lbrack a,b\rbrack$ 上一致收敛于 $h(x)$.
\end{example}

\begin{proof}[解]
当前条件不足。取区间 $[0,1]$，令 $f=0$、$g(x)=x^{-2/3}$，并令
\[
    f_n(x)=n^{1/3}\mathbf1_{(0,1/n)}(x).
\]
则 $g$ 可积，且
\[
    \int_0^1|f_n|^2=n^{-1/3}\to0,
\]
但 $h_n(1)=\int_0^1f_ng=3$，并不趋于零。若补充 $g\in L^2[a,b]$，则由 Cauchy--Schwarz 不等式
\[
    \sup_x|h_n(x)-h(x)|
    \leq\|f_n-f\|_2\|g\|_2\longrightarrow0.
\]
故题面需要核对。
\end{proof}

\begin{example}
设 $f_1(x)$ 在 $\lbrack a,b\rbrack$ 上正常可积,
\[
    f_{n+1}(x)=\int_a^x f_n(t)\,\mathrm{d}t,
    \qquad n=1,2,\ldots.
\]
证明:函数序列 $\{f_n(x)\}$ 在 $\lbrack a,b\rbrack$ 上一致收敛于零.
\end{example}

\begin{proof}[解]
设 $M=\sup_{[a,b]}|f_1|$、$L=b-a$。反复积分可得
\[
    |f_n(x)|\leq M\dfrac{L^{n-1}}{(n-1)!},
    \qquad x\in[a,b].
\]
右端趋于零且与 $x$ 无关，故 $f_n$ 一致收敛于零。
\end{proof}

\begin{example}
\begin{enumerate}[label=(\arabic*)]
    \item 设 $f_n(x)$ 在 $\lbrack a,b\rbrack$ 上可微,且存在 $M>0$,使 $|f_n'(x)|\leq M$.若 $f_n(x)$ 在 $\lbrack a,b\rbrack$ 上逐点收敛于 $f(x)$,则必一致收敛;
    \item 设定义在 $\lbrack 0,1\rbrack$ 上的函数列 $\{f_n(x)\}$ 满足存在 $M>0$ 使得
    $|f_n(x)-f_n(y)|\leq M|x-y|$,且 $\displaystyle\lim\limits_{n\to\infty}f_n(x)=f(x)$,则 $f_n(x)$ 在 $\lbrack 0,1\rbrack$ 上一致收敛于 $f(x)$.
\end{enumerate}
\end{example}

\begin{proof}[解]
两题本质相同。由统一的 Lipschitz 估计，函数族等度连续；逐点极限 $f$ 也满足同一 Lipschitz 估计。任取 $\varepsilon>0$，先把区间分成有限个长度小于 $\varepsilon/(3M)$ 的小区间，再在所有分点处利用逐点收敛统一选取 $N$。任意点与相邻分点比较，即得 $n\geq N$ 时
\[
    \sup_x|f_n(x)-f(x)|<\varepsilon.
\]
故两列均一致收敛。
\end{proof}

\begin{example}
设 $\{f_n(x)\}$ 是 $\lbrack 0,1\rbrack$ 上由递增函数组成的序列,若存在 $f\in C\lbrack 0,1\rbrack$ 使得 $\displaystyle\lim\limits_{n\to\infty}f_n(x)=f(x)$,则 $f_n(x)$ 在 $\lbrack 0,1\rbrack$ 上一致收敛于 $f(x)$.
\end{example}

\begin{proof}[解]
任取 $\varepsilon>0$。由 $f$ 的一致连续性，取分点
$0=x_0<\cdots<x_m=1$，使相邻分点处 $f$ 的振幅小于 $\varepsilon$。在有限个分点上同时取充分大的 $n$，使
$|f_n(x_j)-f(x_j)|<\varepsilon$。若 $x_j\leq x\leq x_{j+1}$，由
$f_n$ 的递增性，
\[
    f_n(x_j)\leq f_n(x)\leq f_n(x_{j+1}),
\]
结合 $f$ 的振幅估计得到 $|f_n(x)-f(x)|<2\varepsilon$，故一致收敛。
\end{proof}

\begin{example}
设 $\{a_n\}$ 是单调收敛于 $0$ 的数列,则
\begin{enumerate}[label=(\arabic*)]
    \item $\displaystyle\sum\limits_{n=1}^{\infty}(-1)^{n-1}a_n\cos nx$ 在 $\lbrack -\pi+\delta,\pi-\delta\rbrack$ $(\delta>0)$ 上一致收敛;
    \item $\displaystyle\sum\limits_{n=1}^{\infty}a_nx\sin nx$ 在 $\lbrack -\pi+\delta,\pi-\delta\rbrack$ $(\delta>0)$ 上一致收敛.
\end{enumerate}
\end{example}

\begin{proof}[解]
\begin{enumerate}[label=(\arabic*)]
    \item $(-1)^{n-1}\cos(nx)$ 的部分和在给定闭区间上一致有界，故由一致 Dirichlet 判别法收敛。
    \item 由三角和公式，
    \[
        \left|x\sum\limits_{k=1}^{N}\sin(kx)\right|
        \leq C_\delta
    \]
    在给定区间上一致成立；再与 $a_n\downarrow0$ 配合使用 Abel 变换，即得一致收敛。
\end{enumerate}
\end{proof}

\begin{example}
\begin{enumerate}[label=(\arabic*)]
    \item 若 $\{na_n\}$ 是单调收敛于 $0$ 的数列,则 $\displaystyle\sum\limits_{n=1}^{\infty}a_n\sin nx$ 在 $(-\infty,+\infty)$ 上一致收敛;
    \item 设 $\{a_n\}$ 是递减正数列.若 $\displaystyle\sum\limits_{n=1}^{\infty}a_n\sin nx$ 在 $\lbrack -a,a\rbrack$ 上一致收敛,则 $na_n\to0$ $(n\to\infty)$.
\end{enumerate}
\end{example}

\begin{proof}[解]
\begin{enumerate}[label=(\arabic*)]
    \item 这是 Chaundy--Jolliffe 判别。令 $m=[1/|x|]$：在 $n\leq m$ 的部分使用 $|\sin(nx)|\leq n|x|$，在 $n>m$ 的部分使用 Abel 变换与三角和估计；两部分均由 $\sup_{k\geq N}k|a_k|$ 控制，故尾和在全实轴上一致趋零。
    \item 由一致 Cauchy 准则，取 $x=1/n$ 并考察 $n\leq k\leq2n$，有
    \[
        \sum\limits_{k=n}^{2n}a_k\sin\dfrac{k}{n}\geq cna_{2n}\longrightarrow0.
    \]
    所以 $na_{2n}\to0$；再由单调性夹逼得到 $na_n\to0$。
\end{enumerate}
\end{proof}

\begin{example}
设 $\displaystyle\sum\limits_{n=1}^{\infty}a_n$ 收敛,且 $\{f_n(x)\}$ 满足
\[
    M\geq f_1(x)\geq f_2(x)\geq\cdots\geq f_n(x)\geq\cdots\geq0.
\]
证明:$\displaystyle\sum\limits_{n=1}^{\infty}a_nf_n(x)$ 在 $\lbrack a,b\rbrack$ 上一致收敛.
\end{example}

\begin{proof}[解]
记 $A_n=\sum_{k=1}^{n}a_k$。对任意 $m>n$ 作 Abel 变换，尾和可写成
\[
    \sum\limits_{k=n}^{m}a_kf_k
    =(A_m-A_{n-1})f_m+
    \sum\limits_{k=n}^{m-1}(A_k-A_{n-1})(f_k-f_{k+1}).
\]
由于 $\sum a_n$ 的尾部分和一致小，而 $f_k$ 的总变差不超过 $M$，右端一致趋于零，故级数一致收敛。
\end{proof}

\begin{example}
设定义在区间 $I$ 上的 $\{f_n(x)\},\{g_n(x)\}$ 满足:
\begin{enumerate}[label=(\arabic*)]
    \item $\displaystyle\sum\limits_{n=1}^{\infty}|f_{n+1}(x)-f_n(x)|$ 在 $I$ 上一致收敛;
    \item $\displaystyle\lim\limits_{n\to\infty}\sup\limits_I|f_n(x)|=0$;
    \item $G_n=\displaystyle\sum\limits_{k=1}^n g_k(x)$ 在 $I$ 上一致有界.
\end{enumerate}
则 $\displaystyle\sum\limits_{n=1}^{\infty}f_n(x)g_n(x)$ 在 $I$ 上一致收敛.
\end{example}

\begin{proof}[解]
对尾和作分部求和：
\[
    \sum\limits_{k=n}^{m}f_kg_k
    =f_mG_m-f_nG_{n-1}
    +\sum\limits_{k=n}^{m-1}(f_k-f_{k+1})G_k.
\]
由 $G_k$ 一致有界、$f_n\rightrightarrows0$，以及
$\sum|f_{n+1}-f_n|$ 的一致收敛性，右端对 $m>n$ 一致趋零。故结论由一致 Cauchy 准则成立。
\end{proof}

\begin{example}
设 $\{f_n(x)\}$ 是区间 $(a,b)$ 内的连续函数序列,并且对任意 $x_0\in(a,b)$,$\{f_n(x_0)\}$ 都是有界的.证明:$\{f_n(x)\}$ 在 $(a,b)$ 的某一非空子区间上一致有界.
\end{example}

\begin{proof}[解]
令
\[
    E_m=\{x\in(a,b):|f_n(x)|\leq m\text{ 对所有 }n\}.
\]
各 $E_m$ 为相对闭集，且由逐点有界性有 $(a,b)=\bigcup_{m=1}^{\infty}E_m$。由 Baire 类定理，某个 $E_m$ 在 $(a,b)$ 中有非空内点；在该开子区间上，全体 $f_n$ 一致有界于 $m$。
\end{proof}

\begin{example}
设 $\{f_n(x)\}$ 在区间 $\lbrack 0,1\rbrack$ 上一致有界,试证存在一子列,其在 $\lbrack 0,1\rbrack$ 的一切有理点上收敛.
\end{example}

\begin{proof}[解]
将 $[0,1]$ 内有理数排列为 $r_1,r_2,\ldots$。由一致有界性，先选子列使其在 $r_1$ 处收敛，再从中选子列使其在 $r_2$ 处收敛，如此继续。取对角子列 $f_{n_k}$，则对每个固定 $j$，它从某项起属于第 $j$ 次所选子列，故在 $r_j$ 处收敛。
\end{proof}

\begin{example}
设 $\mathscr R$ 是区间 $I$ 上的函数族,对任意 $f\in\mathscr R$ 都在 $I$ 上可微,且 $\{f'(x)\mid f\in\mathscr R\}$ 在 $I$ 上一致有界.证明:$\mathscr R$ 在 $I$ 上等度连续.
\end{example}

\begin{proof}[解]
若 $|f'(x)|\leq M$ 对所有 $f\in\mathscr R$、$x\in I$ 成立，则由 Lagrange 中值定理
\[
    |f(x)-f(y)|\leq M|x-y|.
\]
给定 $\varepsilon>0$，统一取 $\delta=\varepsilon/M$ 即得等度连续性。
\end{proof}

\begin{example}
\begin{enumerate}[label=(\arabic*)]
    \item 设 $\{f_n(x)\}$ 是 $\lbrack 0,1\rbrack$ 上二次可导的函数列,且 $f_n(0)=f_n'(0)=0$.若 $|f_n''(x)|\leq1$,则存在子列 $\{f_{n_k}(x)\}$ 在 $\lbrack 0,1\rbrack$ 上一致收敛;
    \item 设 $f_n\in C^{(1)}\lbrack 0,1\rbrack$,且 $|f_n'(x)|\leq x^{-\frac{1}{2}}$ $(0<x\leq1)$,并且 $\displaystyle\int_0^1f_n(x)\,\mathrm{d}x=0$,则存在子列 $\{f_{n_k}(x)\}$ 在 $\lbrack 0,1\rbrack$ 上一致收敛.
\end{enumerate}
\end{example}

\begin{proof}[解]
\begin{enumerate}[label=(\arabic*)]
    \item 由初值与二阶导数界，$|f_n'(x)|\leq1$ 且 $|f_n(x)|\leq1/2$。故函数族一致有界且等度连续，由 Arzel\`a--Ascoli 定理存在一致收敛子列。
    \item 对 $0\leq x<y\leq1$，
    \[
        |f_n(y)-f_n(x)|\leq\int_x^yt^{-1/2}\,\mathrm{d}t
        \leq2\sqrt{y-x},
    \]
    故等度连续。又由积分为零，存在 $c_n\in[0,1]$ 使 $f_n(c_n)=0$，于是 $|f_n(x)|\leq2$。再次应用 Arzel\`a--Ascoli 定理即得一致收敛子列。
\end{enumerate}
\end{proof}

\newpage

\section{函数性质的传递——极限换序}

\subsection{连续性质的传递}

\begin{example}
判别下列级数 $\displaystyle S(x)=\sum\limits_{n=1}^{\infty}f_n(x)$ 在指定区间上的连续性:
\begin{enumerate}[label=(\arabic*)]
    \item $\displaystyle\sum\limits_{n=1}^{\infty}\dfrac{\ln^n x}{n^2}$, $x\in\lbrack e^{-1},e\rbrack$;
    \item $\displaystyle\sum\limits_{n=1}^{\infty}\dfrac{x+n(-1)^n}{n^2+x^2}$, $x\in\lbrack -a,a\rbrack$;
    \item $\displaystyle\sum\limits_{n=1}^{\infty}\dfrac{x^n\sin nx}{n!}$, $x\in\lbrack a,b\rbrack$;
    \item $\displaystyle\sum\limits_{n=1}^{\infty}\left(\dfrac{1}{n}+x\right)^n$, $x\in(-1,1)$.
\end{enumerate}
\end{example}

\begin{proof}[解]
\begin{enumerate}[label=(\arabic*)]
    \item 因 $|\ln x|\leq1$，通项绝对值由 $1/n^2$ 控制，故和函数连续。
    \item 分成
    \[
        \sum\limits_{n=1}^{\infty}\dfrac{x}{n^2+x^2}
        +\sum\limits_{n=1}^{\infty}(-1)^n\dfrac{n}{n^2+x^2}.
    \]
    第一项在紧区间上一致绝对收敛；第二项从 $n>a$ 起系数关于 $n$ 递减且上确界趋零，故一致收敛。和函数连续。
    \item 通项由 $C^n/n!$ 控制，故在任意有限闭区间上一致收敛，和函数连续。
    \item 对任意 $x_0\in(-1,1)$，取 $|x_0|<r<1$，则在 $x_0$ 的充分小邻域内，充分大时 $|x+1/n|\leq q<1$。故级数局部一致收敛，和函数在 $(-1,1)$ 上连续。
\end{enumerate}
\end{proof}

\begin{example}
试求下列极限值:
\begin{enumerate}[label=(\arabic*)]
    \item $\displaystyle\lim\limits_{x\to0^+}\sum\limits_{n=1}^{\infty}\dfrac{1}{2^n n^x}$;
    \item $\displaystyle\lim\limits_{x\to+\infty}\sum\limits_{n=1}^{\infty}\dfrac{x^2}{1+n^2x^2}$;
    \item $\displaystyle\lim\limits_{x\to1^-}\sum\limits_{n=1}^{\infty}\dfrac{(-1)^{n+1}}{n}\dfrac{x^n}{1+x^n}$;
    \item $\displaystyle\lim\limits_{x\to1}\sum\limits_{n=1}^{\infty}\dfrac{x^n\sin\left(\dfrac{n\pi x}{2}\right)}{2^n}$.
\end{enumerate}
\end{example}

\begin{proof}[解]
利用一致收敛或控制收敛逐项取极限：
\begin{enumerate}[label=(\arabic*)]
    \item 由 $2^{-n}$ 控制，极限为 $\sum_{n=1}^{\infty}2^{-n}=\boxed1$。
    \item 因 $0\leq x^2/(1+n^2x^2)\leq1/n^2$，极限为
    $\sum_{n=1}^{\infty}1/n^2=\boxed{\pi^2/6}$。
    \item 一致 Abel 判别允许令 $x\to1^-$，故极限为
    \[
        \dfrac12\sum\limits_{n=1}^{\infty}\dfrac{(-1)^{n+1}}n
        =\boxed{\dfrac12\ln2}.
    \]
    \item 由几何级数控制，令 $x\to1$ 得
    \[
        \sum\limits_{n=1}^{\infty}\dfrac{\sin(n\pi/2)}{2^n}
        =\operatorname{Im}\dfrac{i/2}{1-i/2}
        =\boxed{\dfrac25}.
    \]
\end{enumerate}
\end{proof}

\begin{example}
设 $g_n,h_n\in C\lbrack a,b\rbrack$,且 $g_n(x)\leq g_{n+1}(x)$,$h_n(x)\geq h_{n+1}(x)$,
\[
    \lim\limits_{n\to\infty}g_n(x)=f(x)=\lim\limits_{n\to\infty}h_n(x).
\]
则 $f\in C\lbrack a,b\rbrack$.
\end{example}

\begin{proof}[解]
令 $d_n=h_n-g_n\geq0$。则 $d_n$ 连续、关于 $n$ 单调递减，并逐点趋于零。由 Dini 定理，$d_n\rightrightarrows0$。于是
\[
    0\leq f-g_n\leq d_n,
\]
故 $g_n\rightrightarrows f$。连续函数的一致极限仍连续，因而 $f\in C[a,b]$。
\end{proof}

\begin{example}
设 $f_n,F_n\in C\lbrack a,b\rbrack$ 且 $|f_n(x)|\leq F_n(x)$.若
\[
    \sigma(x)=\sum\limits_{n=1}^{\infty}F_n(x)
\]
在 $\lbrack a,b\rbrack$ 上连续,则 $S(x)=\displaystyle\sum\limits_{n=1}^{\infty}f_n(x)$ 在 $\lbrack a,b\rbrack$ 上连续.
\end{example}

\begin{proof}[解]
由 $|f_n|\leq F_n$ 可知 $F_n\geq0$。连续函数的部分和
$\sigma_N=\sum_{n=1}^{N}F_n$ 单调增加到连续函数 $\sigma$，故由 Dini 定理一致收敛。因此
\[
    \sup_{x\in[a,b]}\sum\limits_{n=N+1}^{\infty}|f_n(x)|
    \leq\sup_{x\in[a,b]}\bigl(\sigma(x)-\sigma_N(x)\bigr)\longrightarrow0.
\]
所以 $\sum f_n$ 一致收敛，其和函数连续。
\end{proof}

\begin{example}
证明:$\displaystyle\sum\limits_{n=-\infty}^{+\infty}\dfrac{1}{(n-x)^2}$ 当 $x$ 不为整数时收敛,周期为 $1$,并且当 $x$ 非整数时和函数连续.
\end{example}

\begin{proof}[解]
在任一不含整数的紧区间 $K$ 上，除有限项外有
\[
    \dfrac1{(n-x)^2}\leq\dfrac4{n^2}\qquad(x\in K),
\]
故双向级数在 $K$ 上一致收敛，和函数在每个非整数点连续。换指标
$m=n-1$ 得
\[
    S(x+1)=\sum\limits_{n\in\mathbb Z}\dfrac1{(n-x-1)^2}
    =\sum\limits_{m\in\mathbb Z}\dfrac1{(m-x)^2}=S(x),
\]
故周期为 $1$。
\end{proof}

\begin{example}[\markstar]
设 $\{x_n\}$ 是 $(0,1)$ 内的一个序列:$0<x_n<1$ 且 $x_i\neq x_j$ $(i\neq j)$.试讨论函数
\[
    f(x)=\sum\limits_{n=1}^{\infty}\dfrac{\operatorname{sgn}(x-x_n)}{2^n}
\]
在 $(0,1)$ 内的连续性.
\end{example}

\begin{proof}[解]
级数由 $2^{-n}$ 控制而一致收敛。固定 $x\notin\{x_n\}$，给定
$\varepsilon>0$，先使尾和小于 $\varepsilon$，再取一个不含有限个首项跳点的邻域，可知 $f$ 在 $x$ 连续。

在 $x=x_k$ 处，第 $k$ 项的左右极限相差 $2/2^k=2^{1-k}$；其余有限个首项在该点附近不变，尾项可一致控制。因此 $f$ 在每个 $x_k$ 处有非零跳跃并不连续。故不连续点集恰为 $\{x_n:n\geq1\}$。
\end{proof}

\subsection{积分性质的传递}

\begin{example}
设 $f\in C^{\infty}(-\infty,+\infty)$,函数列 $\{f^{(n)}(x)\}$ 在任意区间 $\lbrack a,b\rbrack$ 上一致收敛于 $g(x)$,则 $g(x)=Ce^x$.
\end{example}

\begin{proof}[解]
对任意 $a<x<b$，
\[
    f^{(n)}(x)-f^{(n)}(a)=\int_a^xf^{(n+1)}(t)\,\mathrm{d}t.
\]
两列 $f^{(n)}$ 与 $f^{(n+1)}$ 都一致收敛于 $g$，故令 $n\to\infty$ 得
\[
    g(x)-g(a)=\int_a^xg(t)\,\mathrm{d}t.
\]
于是 $g'=g$，故 $g(x)=Ce^x$。
\end{proof}

\begin{example}
证明下列命题:
\begin{enumerate}[label=(\arabic*)]
    \item $\displaystyle\int_0^1\sum\limits_{n=1}^{\infty}\dfrac{x}{n(n+x)}\,\mathrm{d}x=\gamma$,其中 $\gamma$ 为 Euler 常数;
    \item $\displaystyle\int_{\ln2}^{\ln5}\sum\limits_{n=1}^{\infty}ne^{-nx}\,\mathrm{d}x=\dfrac{3}{4}$;
    \item
    \[
        \int_0^{+\infty}\dfrac{x^{a-1}}{1+x}\,\mathrm{d}x
        =\dfrac{1}{a}+\sum\limits_{k=1}^{\infty}(-1)^k
        \left(\dfrac{1}{a+k}-\dfrac{1}{a-k}\right),
        \qquad 0<a<1.
    \]
\end{enumerate}
\end{example}

\begin{proof}[解]
\begin{enumerate}[label=(\arabic*)]
    \item 因
    \[
        \dfrac{x}{n(n+x)}=\dfrac1n-\dfrac1{n+x},
    \]
    逐项积分得到
    \[
        \sum\limits_{n=1}^{\infty}\left[\dfrac1n-\ln\left(1+\dfrac1n\right)\right]=\gamma.
    \]
    \item 逐项积分得
    \[
        \sum\limits_{n=1}^{\infty}(2^{-n}-5^{-n})=1-\dfrac14=\dfrac34.
    \]
    \item 分别在 $(0,1)$ 上展开 $1/(1+x)$，并对 $(1,+\infty)$ 作代换 $t=1/x$ 后展开。逐项积分并合并，得到
    \[
        \dfrac1a+\sum\limits_{k=1}^{\infty}(-1)^k
        \left(\dfrac1{a+k}-\dfrac1{a-k}\right).
    \]
    在 $0<a<1$ 时上述换序可先在 $[\varepsilon,1-\varepsilon]$ 上进行，再令 $\varepsilon\to0^+$。
\end{enumerate}
\end{proof}

\begin{example}
证明下列命题:
\begin{enumerate}[label=(\arabic*)]
    \item $\displaystyle\lim\limits_{n\to\infty}\int_0^1\dfrac{\mathrm{d}x}{1+\left(\dfrac{x}{n}+1\right)^n}=\ln\dfrac{2e}{e+1}$;
    \item $\displaystyle\int_0^1x^{-x}\,\mathrm{d}x=\sum\limits_{n=1}^{\infty}\dfrac{1}{n^n}$.
\end{enumerate}
\end{example}

\begin{proof}[解]
\begin{enumerate}[label=(\arabic*)]
    \item 被积函数逐点趋于 $1/(1+e^x)$ 且由 $1$ 控制，故
    \[
        \lim\limits_{n\to\infty}\int_0^1\dfrac{\mathrm{d}x}{1+(1+x/n)^n}
        =\int_0^1\dfrac{\mathrm{d}x}{1+e^x}
        =\ln\dfrac{2e}{e+1}.
    \]
    \item 展开
    \[
        x^{-x}=e^{-x\ln x}
        =\sum\limits_{k=0}^{\infty}\dfrac{x^k(-\ln x)^k}{k!}.
    \]
    各项非负，可逐项积分；又
    \[
        \int_0^1x^k(-\ln x)^k\,\mathrm{d}x=\dfrac{k!}{(k+1)^{k+1}},
    \]
    故积分等于 $\sum_{n=1}^{\infty}n^{-n}$。
\end{enumerate}
\end{proof}

\begin{example}
设 $f,g\in C(-\infty,+\infty)$,且 $f(x+1)=f(x)$,$g(x+1)=g(x)$.证明:
\[
    \lim\limits_{n\to\infty}\int_0^1f(x)g(nx)\,\mathrm{d}x
    =\left(\int_0^1f(x)\,\mathrm{d}x\right)
     \left(\int_0^1g(x)\,\mathrm{d}x\right).
\]
\end{example}

\begin{proof}[解]
把 $[0,1]$ 分成 $n$ 个小区间并令 $x=(k+t)/n$，得到
\[
    \int_0^1f(x)g(nx)\,\mathrm{d}x
    =\int_0^1g(t)\left[\dfrac1n\sum\limits_{k=0}^{n-1}f\left(\dfrac{k+t}{n}\right)\right]\mathrm{d}t.
\]
方括号内的 Riemann 和关于 $t\in[0,1]$ 一致趋于 $\int_0^1f$，故可取极限，得到题中乘积。
\end{proof}

\begin{example}
设 $h(x),f_n'(x)$ 在 $\lbrack a,b\rbrack$ 上连续,又对 $\lbrack a,b\rbrack$ 中任意 $x_1,x_2$ 和正整数 $n$,有
\[
    |f_n(x_1)-f_n(x_2)|\leq\dfrac{M}{n}|x_1-x_2|,
    \qquad M>0.
\]
证明:
\[
    \lim\limits_{n\to\infty}\int_a^bh(x)f_n'(x)\,\mathrm{d}x=0.
\]
\end{example}

\begin{proof}[解]
由给定 Lipschitz 条件及 $f_n'$ 连续，令 $x_2\to x_1$ 得
\[
    |f_n'(x)|\leq\dfrac Mn.
\]
因此
\[
    \left|\int_a^bh(x)f_n'(x)\,\mathrm{d}x\right|
    \leq\dfrac{M(b-a)}n\max_{[a,b]}|h|\longrightarrow0.
\]
\end{proof}

\begin{example}
设 $g(x),f_n(x)$ 在 $\lbrack a,b\rbrack$ 上有界可积,且对任意 $c\in(a,b)$,$f_n(x)\rightrightarrows0$ 于 $\lbrack c,b\rbrack$ 上,
\[
    \lim\limits_{n\to\infty}\int_a^bf_n(x)\,\mathrm{d}x=1,
    \qquad
    \lim\limits_{x\to a^+}g(x)=A.
\]
试证:
\[
    \lim\limits_{n\to\infty}\int_a^bg(x)f_n(x)\,\mathrm{d}x=A.
\]
\end{example}

\begin{proof}[解]
当前条件不足，因为没有控制 $\int|f_n|$。例如在 $[0,1]$ 上取 $g(x)=x$，可在 $(0,1/n)$ 内取两个符号相反、总积分为零而一阶矩恒为非零的尖峰，再加上积分为 $1$ 的正尖峰；此时仍满足远离零一致趋零及 $\int f_n=1$，但 $\int gf_n$ 不趋于 $g(0)$。

若补充 $f_n\geq0$（或 $\int|f_n|$ 一致有界），则任取 $c>a$，把积分分为 $[a,c]$ 与 $[c,b]$。前一段用 $g(x)\to A$ 控制，后一段用 $f_n\rightrightarrows0$ 控制，再令 $c\to a^+$，即可得到极限 $A$。故原题需要补充条件。
\end{proof}

\begin{example}
设 $\{f_n(x)\}$ 是定义在 $\lbrack -1,1\rbrack$ 上的连续函数序列,在原点附近一致有界.若对任意 $0<\delta<1$,$f_n(x)$ 在 $\lbrack -1,-\delta\rbrack\cup\lbrack \delta,1\rbrack$ 上一致收敛于零,且 $g\in C\lbrack -1,1\rbrack$, $g(0)=0$,证明:
\[
    \lim\limits_{n\to\infty}\int_{-1}^{1}f_n(x)g(x)\,\mathrm{d}x=g(0).
\]
\end{example}

\begin{proof}[解]
结论右端 $g(0)=0$。取 $\delta>0$。在 $|x|<\delta$ 上，$f_n$ 一致有界且
$\sup_{|x|<\delta}|g(x)|\to0$（$\delta\to0$）；在其余闭集上，$f_n$ 一致趋零。因此
\[
    \left|\int_{-1}^{1}f_n(x)g(x)\,\mathrm{d}x\right|
    \leq C\int_{-\delta}^{\delta}|g(x)|\,\mathrm{d}x+o(1).
\]
先令 $n\to\infty$，再令 $\delta\to0$，即得极限为零。
\end{proof}

\subsection{微分性质的传递}

\begin{example}
讨论下列级数 $\displaystyle S(x)=\sum\limits_{n=1}^{\infty}f_n(x)$ 在指定区域的可微性:
\begin{enumerate}[label=(\arabic*)]
    \item $f_n(x)=\arctan\dfrac{x}{n^2}$, $x\in(0,+\infty)$;
    \item $f_n(x)=\dfrac{\sin nx}{n^3}$, $x\in(-\infty,+\infty)$;
    \item $f_n(x)=\sqrt n\tan^n x$, $x\in\left(-\dfrac{\pi}{4},\dfrac{\pi}{4}\right)$.
\end{enumerate}
\end{example}

\begin{proof}[解]
\begin{enumerate}[label=(\arabic*)]
    \item
    \[
        f_n'(x)=\dfrac{n^{-2}}{1+x^2/n^4},\qquad |f_n'(x)|\leq n^{-2}.
    \]
    导数级数在全区间一致收敛，故和函数可微。
    \item $f_n'(x)=\cos(nx)/n^2$，导数级数由 $\sum n^{-2}$ 一致控制，故和函数在实轴上可微。
    \item 在任一紧子区间 $|\tan x|\leq q<1$ 上，原级数及导数级数都由多项式因子乘 $q^n$ 控制，故局部一致收敛并可逐项求导。因此和函数在给定开区间内可微。
\end{enumerate}
\end{proof}

\begin{example}
设 $f(x)$ 在 $(-\infty,+\infty)$ 内有任意阶导数,级数
\[
\begin{aligned}
    \cdots+f^{(n)}(x)+\cdots+f''(x)+f'(x)+f(x)
    &+\int_0^xf(t_1)\,\mathrm{d}t_1\\
    &+\int_0^x\mathrm{d}t_2\int_0^{t_2}f(t_1)\,\mathrm{d}t_1+\cdots
\end{aligned}
\]
按两个方向在 $(-\infty,+\infty)$ 内一致收敛.求该级数的和函数 $F(x)$.
\end{example}

\begin{proof}[解]
按题设可逐项求导，而求导只把双向级数的各项向左平移一位，所以
\[
    F'(x)=F(x).
\]
故 $F(x)=Ce^x$。在 $x=0$ 处所有正重积分项为零，因此
\[
    C=F(0)=\sum\limits_{n=0}^{\infty}f^{(n)}(0).
\]
\end{proof}

\begin{example}
证明:存在 $\xi\in(0,1)$,使得
\[
    \sum\limits_{n=1}^{\infty}\dfrac{\xi^{n-1}}{n+1}=1.
\]
\end{example}

\begin{proof}[解]
令
\[
    F(x)=\sum\limits_{n=1}^{\infty}\dfrac{x^{n-1}}{n+1},\qquad0\leq x<1.
\]
该函数连续且严格递增，并且 $F(0)=1/2<1$，而
$F(x)=[-\ln(1-x)-x]/x^2\to+\infty$（$x\to1^-$）。由介值定理，存在唯一 $\xi\in(0,1)$ 使 $F(\xi)=1$。
\end{proof}

\begin{example}[\markstar]
设 $f_n$ 均在 $\lbrack a,b\rbrack$ 上有原函数,且在 $\lbrack a,b\rbrack$ 上一致收敛于 $f(x)$,则 $f(x)$ 在 $\lbrack a,b\rbrack$ 上也有原函数.
\end{example}

\begin{proof}[解]
取原函数 $F_n$ 并归一化为 $F_n(a)=0$。由 Lagrange 中值定理，
\[
    \|F_n-F_m\|_\infty\leq(b-a)\|f_n-f_m\|_\infty.
\]
故 $F_n$ 一致收敛于某个 $F$。再由导数一致收敛定理，$F'=f$，所以 $f$ 有原函数。
\end{proof}

\begin{example}[\markstar]
设 $f_n$ 在 $\lbrack a,b\rbrack$ 上可积,且在 $\lbrack a,b\rbrack$ 上 $f_n(x)\rightrightarrows f(x)$,则 $f(x)$ 在 $\lbrack a,b\rbrack$ 上可积.
\end{example}

\begin{proof}[解]
任取分割 $P$。由一致收敛，充分大时 $\|f_n-f\|_\infty<\varepsilon$，从而上下 Darboux 和满足
\[
    U(f,P)-L(f,P)
    \leq U(f_n,P)-L(f_n,P)+2\varepsilon(b-a).
\]
先为可积函数 $f_n$ 选择使第一项任意小的分割，再令 $\varepsilon\to0$，即得 $f$ 可积。
\end{proof}

\newpage

\section{连续函数的逼近定理}

\begin{example}
设 $f\in C\lbrack a,b\rbrack$,且对每个非负整数 $n$,满足
\[
    \int_a^bx^nf(x)\,\mathrm{d}x=0.
\]
证明:$f\equiv0$.又若条件改为对大于某个正整数 $n_0$ 的所有 $n$ 成立,则也有相同结论.
\end{example}

\begin{proof}[解]
由 Weierstrass 逼近定理，存在多项式 $p_m$ 一致逼近 $f$。题设给出
$\int_a^bp_m(x)f(x)\,\mathrm{d}x=0$，故令 $m\to\infty$ 得
\[
    \int_a^bf(x)^2\,\mathrm{d}x=0,
\]
从而 $f\equiv0$。

若只对 $n>n_0$ 成立，则对函数 $x^{n_0+1}f(x)$ 的所有多项式矩仍为零，故
$x^{n_0+1}f(x)\equiv0$；由连续性仍有 $f\equiv0$。
\end{proof}

\begin{example}
\begin{enumerate}[label=(\arabic*)]
    \item 设 $f\in C\lbrack -1,1\rbrack$,且对每个非负整数 $n$ 满足 $\displaystyle\int_{-1}^{1}x^{2n}f(x)\,\mathrm{d}x=0$,则 $f$ 必为奇函数;
    \item 设 $f\in C\lbrack -1,1\rbrack$,且对每个非负整数 $n$ 满足 $\displaystyle\int_{-1}^{1}x^{2n+1}f(x)\,\mathrm{d}x=0$,则 $f$ 必为偶函数.
\end{enumerate}
\end{example}

\begin{proof}[解]
\begin{enumerate}[label=(\arabic*)]
    \item 令 $f_e(x)=[f(x)+f(-x)]/2$ 为偶部。题设等价于
    $\int_{-1}^{1}x^{2n}f_e(x)\,\mathrm{d}x=0$。偶连续函数可由 $x^2$ 的多项式一致逼近，故与上一题同理得到 $f_e=0$，即 $f$ 为奇函数。
    \item 对奇部 $f_o(x)=[f(x)-f(-x)]/2$ 作同样论证；$x^{2n+1}f_o(x)$ 为偶函数，最终得到 $f_o=0$，故 $f$ 为偶函数。
\end{enumerate}
\end{proof}

\begin{example}
若实系数多项式序列 $\{p_n(x)\}$ 在 $\mathbb R$ 上一致收敛于函数 $f(x)$,证明:$f(x)$ 必是多项式函数.
\end{example}

\begin{proof}[解]
一致收敛使 $\{p_n\}$ 在 $\mathbb R$ 上一致 Cauchy。故充分大的 $m,n$ 满足
$p_n-p_m$ 在整个实轴上有界。非恒定实多项式在实轴上无界，所以
$p_n-p_m$ 必为常数。固定一个充分大的 $N$，便有 $p_n=p_N+c_n$，且
$c_n$ 收敛。因此 $f=p_N+\lim c_n$，仍为多项式。
\end{proof}

\begin{example}
设 $f\in R\lbrack a,b\rbrack$,则对于每个 $\varepsilon>0$,存在两个多项式 $p(x)$ 和 $P(x)$,使得
\begin{enumerate}[label=(\arabic*)]
    \item $p(x)\leq f(x)\leq P(x)$, $x\in\lbrack a,b\rbrack$;
    \item $\displaystyle\int_a^b[P(x)-p(x)]\,\mathrm{d}x<\varepsilon$.
\end{enumerate}
\end{example}

\begin{proof}[解]
由 Riemann 可积性，可取分割使上下 Darboux 和之差小于 $\varepsilon/3$。把相应上下阶梯函数仅在各分点的充分小邻域内改成连续折线，可得连续函数 $u,v$ 满足
\[
    u\leq f\leq v,qquad \int_a^b(v-u)<\dfrac{2\varepsilon}{3}.
\]
再由 Weierstrass 定理分别用多项式一致逼近 $u$ 与 $v$，并把误差常数向下、向上平移，得到多项式 $p\leq u\leq f\leq v\leq P$，且
$\int(P-p)<\varepsilon$。
\end{proof}

\begin{example}
设 $f\in C\lbrack 0,1\rbrack$,证明:存在每项为奇数项的多项式序列在 $\lbrack 0,1\rbrack$ 上一致收敛于 $f$ 的充分必要条件为 $f(0)=0$.
\end{example}

\begin{proof}[解]
必要性显然，因为所有奇次幂多项式在零点均为零。充分性是
M\"untz--Sz\'asz 定理在指数列 $1,3,5,\ldots$ 上的直接应用：
\[
    \sum\limits_{k=0}^{\infty}\dfrac1{2k+1}=+\infty,
\]
故这些单项式的线性张成在所有满足 $f(0)=0$ 的 $C[0,1]$ 函数中稠密。因此存在只含奇次幂的多项式序列一致逼近 $f$。
\end{proof}

\begin{example}
设 $f\in R\lbrack a,b\rbrack$,且对每个非负整数 $n$ 满足 $\displaystyle\int_a^bx^nf(x)\,\mathrm{d}x=0$,证明:$f$ 在每个连续点上等于 $0$.
\end{example}

\begin{proof}[解]
题设蕴含对任意多项式 $p$ 都有 $\int_a^bp(x)f(x)\,\mathrm{d}x=0$。若某连续点 $c$ 满足 $f(c)>0$，则在其某邻域内 $f$ 有正下界。选取在 $c$ 附近接近 $1$、在该邻域外一致趋于零的非负多项式峰函数 $p_N$，可使
\[
    \int_a^bp_N(x)f(x)\,\mathrm{d}x>0,
\]
矛盾；$f(c)<0$ 同理。因此每个连续点处都有 $f(c)=0$。
\end{proof}

\begin{example}
设 $f\in C\lbrack 0,1\rbrack$,若
\[
    \int_0^1f\left(x^{\frac{1}{2n+1}}\right)\,\mathrm{d}x=0,
\]
则 $f(x)\equiv0$.
\end{example}

\begin{proof}[解]
对每个非负整数 $n$ 作代换 $x=t^{2n+1}$，题设成为
\[
    (2n+1)\int_0^1t^{2n}f(t)\,\mathrm{d}t=0.
\]
而 $1,t^2,t^4,\ldots$ 的线性组合在 $C[0,1]$ 中稠密，因为任意
$h(t)$ 可写成连续函数 $H(t^2)$。故可用这些多项式一致逼近 $f$，从而
$\int_0^1f^2=0$，即 $f\equiv0$。
\end{proof}

\begin{example}
\begin{enumerate}[label=(\arabic*)]
    \item 设 $f$ 在 $\lbrack a,b\rbrack$ 上有界且有原函数,$g\in C\lbrack a,b\rbrack$,则 $f(x)g(x)$ 在 $\lbrack a,b\rbrack$ 上有原函数;
    \item 设 $f(x)$ 在 $\lbrack a,b\rbrack$ 上有界且有原函数,$g:\lbrack 0,1\rbrack\to\lbrack a,b\rbrack$ 有连续导数且 $g'(x)>0$,则 $f[g(x)]$ 在 $\lbrack 0,1\rbrack$ 上有原函数.
\end{enumerate}
\end{example}

\begin{proof}[解]
\begin{enumerate}[label=(\arabic*)]
    \item 设 $F'=f$。用多项式 $q_m$ 一致逼近 $g$。对每个多项式 $q$，
    \[
        fq=(Fq)'-Fq'
    \]
    有原函数；又 $fq_m\rightrightarrows fg$，由前述一致极限定理，$fg$ 也有原函数。
    \item $g'>0$，故 $1/g'$ 连续。函数
    \[
        f(g(x))g'(x)=\dfrac{\mathrm{d}}{\mathrm{d}x}F(g(x))
    \]
    有原函数。再对它与连续函数 $1/g'$ 应用 (1)，即得 $f(g(x))$ 有原函数。
\end{enumerate}
\end{proof}

\newpage

\section{幂级数}

\subsection{收敛半径与收敛域}

\begin{example}
试求下列幂级数的收敛半径 $R$:
\begin{enumerate}[label=(\arabic*)]
    \item $\displaystyle\sum\limits_{n=1}^{\infty}\dfrac{(n!)^2}{(2n)!}x^n$;
    \item $\displaystyle\sum\limits_{n=1}^{\infty}\ln\dfrac{3n-2}{3n+2}x^n$;
    \item $\displaystyle\sum\limits_{n=1}^{\infty}\ln\left(\cos\dfrac{1}{3n}\right)x^n$;
    \item $\displaystyle\sum\limits_{n=1}^{\infty}\dfrac{1}{n!}\left(\dfrac{nx}{e}\right)^n$;
    \item $\displaystyle\sum\limits_{n=1}^{\infty}\dfrac{n^n}{n!}e^{-n\alpha}x^n$;
    \item $\displaystyle\sum\limits_{n=1}^{\infty}\dfrac{x^n}{a^n+b^n}$.
\end{enumerate}
\end{example}

\begin{proof}[解]
% TODO: 补充本题解答。
\end{proof}

\begin{example}
试求下列幂级数的收敛半径 $R$:
\begin{enumerate}[label=(\arabic*)]
    \item $\displaystyle\sum\limits_{n=0}^{\infty}5^n x^{3n}$;
    \item $\displaystyle\sum\limits_{n=1}^{\infty}n^n x^{n^2}$;
    \item $\displaystyle\sum\limits_{n=0}^{\infty}n!x^{n!}$.
\end{enumerate}
\end{example}

\begin{proof}[解]
% TODO: 补充本题解答。
\end{proof}

\begin{example}
若存在 $\alpha>0,l>0$,使得
\[
    \lim\limits_{n\to\infty}|a_nn^\alpha|=l,
\]
则 $\displaystyle\sum\limits_{n=0}^{\infty}a_nx^n$ 的收敛半径 $R=1$.
\end{example}

\begin{proof}[解]
% TODO: 补充本题解答。
\end{proof}

\begin{example}
设 $a_n>0$, $S_n=\displaystyle\sum\limits_{k=0}^{n}a_k$.若
\[
    \lim\limits_{n\to\infty}S_n=+\infty,
    \qquad
    \lim\limits_{n\to\infty}\dfrac{a_n}{S_n}=0,
\]
则 $\displaystyle\sum\limits_{n=0}^{\infty}a_nx^n$ 的收敛半径 $R=1$.
\end{example}

\begin{proof}[解]
% TODO: 补充本题解答。
\end{proof}

\begin{example}
设 $a_n>0$, $\displaystyle\sum\limits_{n=0}^{\infty}a_nx^n$ 的收敛半径是 $1$,则
\[
    I=\sum\limits_{n=0}^{\infty}S_nx^n,
    \qquad
    S_n=\sum\limits_{k=0}^{n}a_k,
\]
的收敛半径也是 $1$.
\end{example}

\begin{proof}[解]
% TODO: 补充本题解答。
\end{proof}

\begin{example}
证明:
\[
    \sum\limits_{\substack{k+j=n\\0\leq k,j\leq n}}C_{2k}^{k}C_{2j}^{j}=4^n,
    \qquad n=0,1,2,\ldots.
\]
\end{example}

\begin{proof}[解]
% TODO: 补充本题解答。
\end{proof}

\begin{example}
试求下列幂级数的收敛域:
\begin{enumerate}[label=(\arabic*)]
    \item $\displaystyle\sum\limits_{n=1}^{\infty}\dfrac{x^n}{1+\dfrac{1}{2}+\cdots+\dfrac{1}{n}}$;
    \item $\displaystyle\sum\limits_{n=0}^{\infty}C_{\alpha}^{n}x^n$.
\end{enumerate}
\end{example}

\begin{proof}[解]
% TODO: 补充本题解答。
\end{proof}

\begin{example}
试求下列幂级数的收敛区域:
\begin{enumerate}[label=(\arabic*)]
    \item $\displaystyle\sum\limits_{n=1}^{\infty}(-1)^{n-1}\dfrac{2n+3}{3n^2+4}x^{2n+1}$;
    \item $\displaystyle\sum\limits_{n=1}^{\infty}(\sqrt[n]{a}-1)x^n$, $a>0$;
    \item $\displaystyle\sum\limits_{n=0}^{\infty}\sqrt{2-2\cos\dfrac{\pi}{2n}}\,x^n$;
    \item $\displaystyle\sum\limits_{n=1}^{\infty}\dfrac{x^n}{n^p\ln n}$, $p>0$;
    \item $\displaystyle\sum\limits_{n=1}^{\infty}\dfrac{[2+(-1)^n]^n}{n}x^n$;
    \item $\displaystyle\sum\limits_{n=1}^{\infty}\dfrac{(1+\dfrac{1}{n})(-1)^n}{n^2}x^n$.
\end{enumerate}
\end{example}

\begin{proof}[解]
% TODO: 补充本题解答。
\end{proof}

\begin{example}
试用变量替换,并借助幂级数法求下列函数项级数的收敛域:
\begin{enumerate}[label=(\arabic*)]
    \item $\displaystyle\sum\limits_{n=1}^{\infty}\dfrac{\sqrt[3]{2n+1}-\sqrt[3]{2n-1}}{\sqrt n}(x+3)^n$;
    \item $\displaystyle\sum\limits_{n=1}^{\infty}\dfrac{n+2}{2^n}e^{-nx}$;
    \item $\displaystyle\sum\limits_{n=1}^{\infty}\dfrac{1}{3+2n}\left(\dfrac{1-x}{1+x}\right)^n$;
    \item $\displaystyle\sum\limits_{n=1}^{\infty}\sqrt n(\tan x)^n$;
    \item $\displaystyle\sum\limits_{n=1}^{\infty}\dfrac{n4^n}{3^n}\dfrac{x^n}{(1-x)^n}$;
    \item $\displaystyle\sum\limits_{n=1}^{\infty}\dfrac{2^{-n}}{\sqrt n}x^n$.
\end{enumerate}
\end{example}

\begin{proof}[解]
% TODO: 补充本题解答。
\end{proof}

\begin{example}
设 $a_0=0,a_1=1,a_{n+1}=a_{n-1}+a_n$,则 $\displaystyle\sum\limits_{n=0}^{\infty}a_nx^n$ 在 $|x|<\dfrac{1}{2}$ 时收敛.
\end{example}

\begin{proof}[解]
% TODO: 补充本题解答。
\end{proof}

\begin{example}
若 $\displaystyle\sum\limits_{n=0}^{\infty}a_n^2x^n$ 的收敛域为 $\lbrack -1,1\rbrack$,则 $\displaystyle\sum\limits_{n=0}^{\infty}\dfrac{a_n}{n}x^n$ 的收敛域为 $\lbrack -1,1\rbrack$.
\end{example}

\begin{proof}[解]
% TODO: 补充本题解答。
\end{proof}

\begin{example}
若 $\displaystyle\sum\limits_{n=1}^{\infty}a_n$ 在 Cesaro 意义下可求和,则 $\displaystyle\sum\limits_{n=1}^{\infty}a_nx^n$ 在 $(-1,1)$ 收敛.
\end{example}

\begin{proof}[解]
% TODO: 补充本题解答。
\end{proof}

\subsection{幂级数求和}

\begin{example}
设 $a_n>0$.若
\[
    \lim\limits_{x\to1^-}\sum\limits_{n=0}^{\infty}a_nx^n=S,
\]
则 $\displaystyle\sum\limits_{n=0}^{\infty}a_n=S$.
\end{example}

\begin{proof}[解]
% TODO: 补充本题解答。
\end{proof}

\begin{example}[Tauber 定理的推广 $*$]
设 $f(x)=\displaystyle\sum\limits_{n=0}^{\infty}a_nx^n$ 的收敛半径为 $1$.若
\[
    \lim\limits_{n\to\infty}\dfrac{1}{n}\sum\limits_{k=1}^{n}ka_k=0,
    \qquad
    \lim\limits_{x\to1^-}f(x)=l,
\]
则 $\displaystyle\sum\limits_{n=0}^{\infty}a_n=l$.
\end{example}

\begin{proof}[解]
% TODO: 补充本题解答。
\end{proof}

\begin{example}
求下列幂级数的和 $S(x)$:
\begin{enumerate}[label=(\arabic*)]
    \item $\displaystyle\sum\limits_{n=1}^{\infty}nx^n$, $|x|<1$;
    \item $\displaystyle\sum\limits_{n=1}^{\infty}\dfrac{x^n}{n}$, $-1\leq x<1$;
    \item $\displaystyle\sum\limits_{n=1}^{\infty}\dfrac{x^n}{n(n+1)(n+2)}$, $|x|\leq1$;
    \item $\displaystyle\sum\limits_{n=1}^{\infty}\dfrac{x^{2n-1}}{2n-1}$, $|x|<1$.
\end{enumerate}
\end{example}

\begin{proof}[解]
% TODO: 补充本题解答。
\end{proof}

\begin{example}
求下列幂级数的和 $S(x)$:
\begin{enumerate}[label=(\arabic*)]
    \item $1+2x-4x^3-5x^4+7x^6+8x^7-\cdots$, $|x|<1$;
    \item $\displaystyle\sum\limits_{n=1}^{\infty}n^3x^n$, $|x|<1$;
    \item $\displaystyle\sum\limits_{n=1}^{\infty}\dfrac{(-1)^{n-1}x^{2n}}{n(2n-1)}$, $|x|<1$;
    \item 设 $\displaystyle S(x)=\sum\limits_{n=1}^{\infty}\dfrac{x^n}{n^2}$,求 $S(x)+S(1-x)$, $0<x<1$.
\end{enumerate}
\end{example}

\begin{proof}[解]
% TODO: 补充本题解答。
\end{proof}

\begin{example}
试用幂级数法求下列常数项级数的和 $S$:
\begin{enumerate}[label=(\arabic*)]
    \item $\displaystyle S=\sum\limits_{n=2}^{\infty}\dfrac{(-1)^n}{n^2+n-2}$;
    \item $\displaystyle S=\sum\limits_{n=1}^{\infty}\dfrac{1}{n(2n+1)2^n}$;
    \item $\displaystyle S=\sum\limits_{n=0}^{\infty}\dfrac{(-1)^n}{3n+1}$;
    \item $\displaystyle S=1-\dfrac{1}{5}+\dfrac{1}{7}-\dfrac{1}{11}+\cdots+\dfrac{1}{6n-5}-\dfrac{1}{6n-1}+\cdots$.
\end{enumerate}
\end{example}

\begin{proof}[解]
% TODO: 补充本题解答。
\end{proof}

\begin{example}
试用公式
\[
    \dfrac{1}{k(k+1)\cdots(k+n)}
    =\dfrac{1}{n!}\int_0^1(1-x)^nx^{k-1}\,\mathrm{d}x
\]
求下列级数的和:
\begin{enumerate}[label=(\arabic*)]
    \item $\displaystyle\dfrac{1}{1\cdot2\cdot3}+\dfrac{1}{3\cdot4\cdot5}+\dfrac{1}{5\cdot6\cdot7}+\cdots$;
    \item $\displaystyle\dfrac{1}{1\cdot2\cdot3\cdot4}+\dfrac{1}{4\cdot5\cdot6\cdot7}+\dfrac{1}{7\cdot8\cdot9\cdot10}+\cdots$.
\end{enumerate}
\end{example}

\begin{proof}[解]
% TODO: 补充本题解答。
\end{proof}

\begin{example}
计算极限
\[
    \lim\limits_{\substack{n\to\infty\\m\to\infty}}
    \sum\limits_{i=1}^{m}\sum\limits_{j=1}^{n}\dfrac{(-1)^{i+j}}{i+j}.
\]
\end{example}

\begin{proof}[解]
% TODO: 补充本题解答。
\end{proof}

\begin{example}
设 $\displaystyle\lim\limits_{n\to\infty}a_n=l$,则 $\displaystyle S(x)=\sum\limits_{n=1}^{\infty}a_nx^n$ 在 $(-1,1)$ 上收敛,且
\[
    \lim\limits_{x\to1^-}(1-x)S(x)=l.
\]
\end{example}

\begin{proof}[解]
% TODO: 补充本题解答。
\end{proof}

\begin{example}
设 $a_n>0,b_n>0$,且
\[
    S_1(x)=\sum\limits_{n=0}^{\infty}a_nx^n,
    \qquad
    S_2(x)=\sum\limits_{n=0}^{\infty}b_nx^n
\]
的收敛半径均为 $1$.如果
\[
    \lim\limits_{x\to1^-}S_1(x)=\lim\limits_{x\to1^-}S_2(x)=+\infty,
    \qquad
    \lim\limits_{n\to\infty}\dfrac{a_n}{b_n}=l,
\]
则
\[
    \lim\limits_{x\to1^-}\dfrac{S_1(x)}{S_2(x)}=l.
\]
\end{example}

\begin{proof}[解]
% TODO: 补充本题解答。
\end{proof}

\begin{example}[\markstar]
设幂级数 $\displaystyle\sum\limits_{n=0}^{\infty}a_nx^n$ 的系数满足
\[
    a_n+Aa_{n-1}+Ba_{n-2}=0,
\]
且此幂级数在点 $x=x_0$ 处收敛.证明:
\[
    S=\sum\limits_{n=0}^{\infty}a_nx_0^n
    =\dfrac{a_0+(a_1+Aa_0)x_0}{1+Ax_0+Bx_0^2}.
\]
\end{example}

\begin{proof}[解]
% TODO: 补充本题解答。
\end{proof}

\begin{example}[\markstar\ 母函数法]
设 $a_0=1,a_1=0,a_2=-5$,
\[
    a_n=4a_{n-1}-5a_{n-2}+2a_{n-3}.
\]
证明:
\[
    a_n=-2^{n+2}+3n+5.
\]
\end{example}

\begin{proof}[解]
% TODO: 补充本题解答。
\end{proof}

\subsection{幂级数展开——Taylor 级数}

\begin{example}
求下列函数在 $x=0$ 处的 Taylor 级数展式:
\begin{enumerate}[label=(\arabic*)]
    \item $f(x)=\dfrac{3x+8}{(2x-3)(x^2+4)}$;
    \item $f(x)=\ln(4+3x-x^2)$;
    \item $f(x)=\sin^4x$.
\end{enumerate}
\end{example}

\begin{proof}[解]
% TODO: 补充本题解答。
\end{proof}

\begin{example}
求下列函数的 Maclaurin 展开式:
\begin{enumerate}[label=(\arabic*)]
    \item $f(x)=\ln(1+x+x^2+x^3)$;
    \item $f(x)=\ln^2(1+x)$;
    \item $f(x)=\dfrac{\ln(1+x)}{1+x}$.
\end{enumerate}
\end{example}

\begin{proof}[解]
% TODO: 补充本题解答。
\end{proof}

\begin{example}
利用待定系数法求下列函数的 Maclaurin 级数展开式:
\begin{enumerate}[label=(\arabic*)]
    \item $f(x)=\dfrac{x\sin\alpha}{1-2x\cos\alpha+x^2}$;
    \item $f(x)=\dfrac{\arcsin x}{\sqrt{1-x^2}}$;
    \item $f(x)=\dfrac{1}{\sqrt{1-2ax+x^2}}$.
\end{enumerate}
\end{example}

\begin{proof}[解]
% TODO: 补充本题解答。
\end{proof}

\begin{example}
利用微分积分法求下列函数的 Maclaurin 级数展开式:
\begin{enumerate}[label=(\arabic*)]
    \item $f(x)=\arctan\dfrac{x+3}{x-3}$;
    \item $f(x)=\ln\left(x+\sqrt{x^2+1}\right)$;
    \item $f(x)=\arccos(1-2x^2)$.
\end{enumerate}
\end{example}

\begin{proof}[解]
% TODO: 补充本题解答。
\end{proof}

\subsection{幂级数展开式的应用}

\methodsection{(A) 求级数的和}

\begin{example}
试求下列级数的和:
\begin{enumerate}[label=(\arabic*)]
    \item $\displaystyle S(x)=\sum\limits_{n=0}^{\infty}\dfrac{\ln^n x}{n!}$, $x>0$;
    \item $\displaystyle S(x)=\sum\limits_{n=0}^{\infty}\dfrac{(-1)^n\ln^n x}{2^n n!}$;
    \item $\displaystyle S(x)=\sum\limits_{n=0}^{\infty}\dfrac{2^n(n+1)}{n!}$.
\end{enumerate}
\end{example}

\begin{proof}[解]
% TODO: 补充本题解答。
\end{proof}

\begin{example}
设 $\dfrac{1}{1-x-x^2}$ 的 Maclaurin 级数为 $\displaystyle\sum\limits_{n=0}^{\infty}a_nx^n$,试证明:
\[
    \sum\limits_{n=0}^{\infty}\dfrac{a_{n+1}}{a_na_{n+2}}=2.
\]
\end{example}

\begin{proof}[解]
% TODO: 补充本题解答。
\end{proof}

\methodsection{(B) 求积分值}

\begin{example}
求下列积分值:
\begin{enumerate}[label=(\arabic*)]
    \item $\displaystyle\int_0^1e^x\ln x\,\mathrm{d}x$;
    \item $\displaystyle\int_0^{+\infty}\dfrac{x}{e^{2\pi x}-1}\,\mathrm{d}x$;
    \item $\displaystyle\int_0^x\dfrac{\sin t}{t}\,\mathrm{d}t$;
    \item $\displaystyle\int_0^1\ln x\ln(1-x)\,\mathrm{d}x$.
\end{enumerate}
\end{example}

\begin{proof}[解]
% TODO: 补充本题解答。
\end{proof}

\begin{example}
证明下列等式:
\begin{enumerate}[label=(\arabic*)]
    \item $\displaystyle\int_0^1\dfrac{\ln(1+x)}{x}\,\mathrm{d}x=\dfrac{\pi^2}{12}$;
    \item $\displaystyle\int_0^{\frac{\pi}{2}}\dfrac{\mathrm{d}\theta}{\sqrt{1-k^2\sin^2\theta}}$, $0\leq k<1$.
\end{enumerate}
\end{example}

\begin{proof}[解]
% TODO: 补充本题解答。
\end{proof}

\methodsection{(C) 求导数 \texorpdfstring{$f^{(n)}(0)$}{f(n)(0)}}

\begin{example}
设
\[
    f(x)=\dfrac{1-x+x^2}{1+x+x^2},
\]
试求 $f^{(n)}(0)$.
\end{example}

\begin{proof}[解]
% TODO: 补充本题解答。
\end{proof}

\newpage

\section{Fourier 级数}

\subsection{Fourier 系数}

\begin{example}
设
\[
    \dfrac{a_0}{2}+\sum\limits_{k=1}^{\infty}(a_k\cos kx+b_k\sin kx)
\]
在 $\lbrack -\pi,\pi\rbrack$ 上一致收敛,试证它必是 $\lbrack -\pi,\pi\rbrack$ 上其和函数的 Fourier 级数.
\end{example}

\begin{proof}[解]
% TODO: 补充本题解答。
\end{proof}

\begin{example}[\markstar]
三角多项式
\[
    \dfrac{a_0}{2}+\sum\limits_{k=1}^{n}(a_k\cos kx+b_k\sin kx)
\]
的 Fourier 级数是它本身.
\end{example}

\begin{proof}[解]
% TODO: 补充本题解答。
\end{proof}

\begin{example}
设 $f(x)=\pi-x$, $x\in(0,\pi)$.
\begin{enumerate}[label=(\arabic*)]
    \item 将 $f(x)$ 展开为正弦级数;
    \item 判断该级数在 $(0,\pi)$ 内是否一致收敛.
\end{enumerate}
\end{example}

\begin{proof}[解]
% TODO: 补充本题解答。
\end{proof}

\begin{example}
试将
\[
    f(x)=
    \begin{cases}
        1-x,&0\leq x\leq2,\\
        x-3,&2<x\leq4
    \end{cases}
\]
在 $\lbrack 0,4\rbrack$ 上展开为余弦级数,并证明
\[
    \sum\limits_{n=0}^{\infty}\dfrac{1}{(2n+1)^2}=\dfrac{\pi^2}{8}.
\]
\end{example}

\begin{proof}[解]
% TODO: 补充本题解答。
\end{proof}

\begin{example}
求
\[
    f(x)=
    \begin{cases}
        x,&0\leq x<1,\\
        2-x,&1\leq x\leq2
    \end{cases}
\]
在 $\lbrack 0,2\rbrack$ 上的 Fourier 展开式.
\end{example}

\begin{proof}[解]
% TODO: 补充本题解答。
\end{proof}

\begin{example}
求
\[
    f(x)=
    \begin{cases}
        0,&-\pi\leq x\leq0,\\
        \sin x,&0<x\leq\pi
    \end{cases}
\]
在 $\lbrack -\pi,\pi\rbrack$ 上的 Fourier 展开式.
\end{example}

\begin{proof}[解]
% TODO: 补充本题解答。
\end{proof}

\begin{example}
设 $f(x)=x$, $x\in\lbrack 0,\dfrac{\pi}{2}\rbrack$,试将 $f(x)$ 展开为
\[
    \sum\limits_{n=1}^{\infty}b_{2n-1}\sin(2n-1)x
\]
型的三角级数.
\end{example}

\begin{proof}[解]
% TODO: 补充本题解答。
\end{proof}

\begin{example}
已知 $f(x)=x$ 在 $(-\pi,\pi)$ 上的 Fourier 级数展开式为
\[
    x=\sum\limits_{n=1}^{\infty}\dfrac{2(-1)^{n-1}}{n}\sin nx,
    \qquad |x|<\pi.
\]
试求函数 $\varphi(x)=x\sin x$ 在 $\lbrack -\pi,\pi\rbrack$ 上的 Fourier 展开式.
\end{example}

\begin{proof}[解]
% TODO: 补充本题解答。
\end{proof}

\begin{example}
设 $f(x)$ 有 Fourier 系数
\[
    f(x)\sim\dfrac{a_0}{2}+\sum\limits_{n=1}^{\infty}(a_n\cos nx+b_n\sin nx).
\]
试证逐项乘 $\sin x$ 可得 $f(x)\sin x$ 的 Fourier 级数.
\end{example}

\begin{proof}[解]
% TODO: 补充本题解答。
\end{proof}

\begin{example}[\markstar]
\begin{enumerate}[label=(\arabic*)]
    \item 设 $f(x)$ 以 $2\pi$ 为周期,在 $(0,2\pi)$ 内有界.若 $f(x)$ 单调递减,试证 $b_n\geq0$;若单调递增,则 $b_n\leq0$;
    \item 设 $f'(x)$ 在 $(0,2\pi)$ 内有界.若 $f'(x)$ 单调递增,试证 $a_n\geq0$;若单调递减,则 $a_n\leq0$;
    \item 设 $f(x)$ 在 $(0,2\pi)$ 内可积.若 $F(x)=\displaystyle\int_0^x\left(f(t)-\dfrac{a_0}{2}\right)\mathrm{d}t$ 单调递减,试证 $a_n\geq0$;若 $F$ 单调递增,则 $a_n\leq0$.
\end{enumerate}
\end{example}

\begin{proof}[解]
% TODO: 补充本题解答。
\end{proof}

\begin{example}
设 $f$ 是以 $2\pi$ 为周期的函数且存在 $\alpha\in(0,1\rbrack$,使 $f$ 满足 $\alpha$ 阶 Lipschitz 条件
\[
    |f(x)-f(y)|\leq L|x-y|^\alpha.
\]
证明:
\[
    a_n=o\left(\dfrac{1}{n^\alpha}\right),
    \qquad
    b_n=o\left(\dfrac{1}{n^\alpha}\right).
\]
\end{example}

\begin{proof}[解]
% TODO: 补充本题解答。
\end{proof}

\subsection{收敛定理}

\begin{example}
证明:若 $f(x)$ 是周期为 $2\pi$ 的连续函数,在 $\lbrack -\pi,\pi\rbrack$ 上分段光滑,则 $f(x)$ 的 Fourier 级数一致收敛于 $f(x)$.
\end{example}

\begin{proof}[解]
% TODO: 补充本题解答。
\end{proof}

\begin{example}
构造两个以 $2\pi$ 为周期的连续函数,使得其 Fourier 级数在 $\lbrack 0,\pi\rbrack$ 上一致收敛于 $0$.
\end{example}

\begin{proof}[解]
% TODO: 补充本题解答。
\end{proof}

\begin{example}
设 $f\in C\lbrack 0,1\rbrack$ 且 $f(0)=0$.若
\[
    \int_0^1f(x)\sin(n\pi x)\,\mathrm{d}x=0,
\]
则 $f\equiv0$.
\end{example}

\begin{proof}[解]
% TODO: 补充本题解答。
\end{proof}

\begin{example}
设 $f(x)$ 是以 $2\pi$ 为周期的函数,在 $\lbrack -\pi,\pi\rbrack$ 上可积.已知其 Fourier 级数部分和可表示为 Dirichlet 积分
\[
    S_n(x)=\dfrac{1}{\pi}\int_{-\pi}^{\pi}f(x+t)
    \dfrac{\sin\left((n+\dfrac{1}{2})t\right)}{2\sin\dfrac{t}{2}}\,\mathrm{d}t,
\]
其中
\[
    D_n(t)=\dfrac{\sin\left((n+\dfrac{1}{2})t\right)}{2\sin\dfrac{t}{2}}
    =\dfrac{1}{2}+\cos t+\cos2t+\cdots+\cos nt.
\]
又令 $\displaystyle\sigma_n(x)=\dfrac{1}{n}\sum\limits_{k=0}^{n-1}S_k(x)$.试证:
\begin{enumerate}[label=(\arabic*)]
    \item $\displaystyle\sum\limits_{k=0}^{n-1}D_k(x)=\dfrac{1}{2}\left(\dfrac{\sin\dfrac{nx}{2}}{\sin\dfrac{x}{2}}\right)^2$;
    \item $\displaystyle\dfrac{1}{2n\pi}\int_{-\pi}^{\pi}\left(\dfrac{\sin\dfrac{nx}{2}}{\sin\dfrac{x}{2}}\right)^2\mathrm{d}x=1$;
    \item 对任意 $\delta>0$,
    \[
        \dfrac{1}{n\pi}\int_{\delta}^{\pi}\left(\dfrac{\sin\dfrac{nx}{2}}{\sin\dfrac{x}{2}}\right)^2\mathrm{d}x\to0,
        \qquad n\to\infty;
    \]
    \item 若 $f(x)$ 是以 $2\pi$ 为周期的连续函数,则 $\sigma_n(x)\rightrightarrows f(x)$ 于 $\lbrack -\pi,\pi\rbrack$ 上.
\end{enumerate}
\end{example}

\begin{proof}[解]
% TODO: 补充本题解答。
\end{proof}

\begin{example}
求
\[
    f(x)=\ln(1+2r\cos x+r^2),
    \qquad |r|<1,
\]
的 Fourier 级数展开式.
\end{example}

\begin{proof}[解]
% TODO: 补充本题解答。
\end{proof}

\begin{example}
已知
\[
    x=2\sum\limits_{n=1}^{\infty}\dfrac{(-1)^{n+1}\sin nx}{n},
    \qquad -\pi<x<\pi.
\]
试求函数 $x^2,x^3$ 的 Fourier 级数展开式.
\end{example}

\begin{proof}[解]
% TODO: 补充本题解答。
\end{proof}

\begin{example}
\begin{enumerate}[label=(\arabic*)]
    \item 证明 $\displaystyle\sum\limits_{n=2}^{\infty}\dfrac{1}{\ln n}\sin nx$ 不可能为某函数的 Fourier 级数;
    \item 证明 $\displaystyle\sum\limits_{n=9}^{\infty}\dfrac{\sin nx}{\ln\ln n}$ 不可能为某函数的 Fourier 级数.
\end{enumerate}
\end{example}

\begin{proof}[解]
% TODO: 补充本题解答。
\end{proof}

\subsection{Parseval 等式}

\begin{example}
求函数
\[
    f(x)=\left(x-\dfrac{\pi}{2}\right)^2,
    \qquad x\in\lbrack 0,2\pi\rbrack,
\]
的 Fourier 级数,并求级数 $\displaystyle\sum\limits_{n=1}^{\infty}\dfrac{1}{n^4}$ 的和.
\end{example}

\begin{proof}[解]
% TODO: 补充本题解答。
\end{proof}

\begin{example}
\begin{enumerate}[label=(\arabic*)]
    \item 设 $f(x)$ 在 $\lbrack 0,\pi\rbrack$ 上有连续导数,且 $f'(x)$ 在 $\lbrack 0,\pi\rbrack$ 上分段光滑,$\displaystyle\int_0^\pi f(x)\,\mathrm{d}x=0$.证明相应 Wirtinger 型不等式;
    \item 设 $f(x)$ 在 $\mathbb R$ 上有连续导数,周期为 $l>0$,且 $\displaystyle\int_0^lf(x)\,\mathrm{d}x=0$.证明
    \[
        \int_0^l|f'(x)|^2\,\mathrm{d}x
        \geq\dfrac{4\pi^2}{l^2}\int_0^l|f(x)|^2\,\mathrm{d}x.
    \]
    等号成立当且仅当
    \[
        f(x)=a_1\cos\dfrac{2\pi x}{l}+b_1\sin\dfrac{2\pi x}{l}.
    \]
\end{enumerate}
\end{example}

\begin{proof}[解]
% TODO: 补充本题解答。
\end{proof}

\end{document}

